% APA 7: márgenes de 1 pulgada (2.54 cm) en los cuatro lados
\usepackage[
top=2.54cm,     % Margen superior
bottom=2.54cm,  % Margen inferior
left=2.54cm,    % Margen izquierdo
right=2.54cm,   % Margen derecho
% showframe      % Opcional: muestra un marco alrededor del área de texto y los márgenes para depuración
]{geometry}

%Automatic indexes in the table of contents
\usepackage[nottoc,notlot]{tocbibind}
%	\renewcommand{\tableofcontents}{\begingroup
%		\tocsection
%		\tocchapter
%		\tocfile{\contentsname}{toc}
%	\endgroup}
%	\renewcommand{\listoffigures}{\begingroup
%    	\tocsection
%        \tocchapter
%    	\tocfile{\listfigurename}{lof}
% 	\endgroup}
%	\renewcommand{\listoftables}{\begingroup
%    	\tocsection
%    	\tocfile{\listtablename}{lot}
% 	\endgroup}

\usepackage[utf8]{inputenc} % coding
\usepackage{acronym}
%Verwendete Schrift z.B. lmodern
%\usepackage{lmodern}
%\usepackage{libertine}
%\usepackage{mathpazo,eulervm}
%\usepackage{tgpagella,eulervm}
%\usepackage[charter]{mathdesign}
%\usepackage[osf,sc]{mathpazo}  % APA 7: desactivado (Palatino con cifras elzevirianas)
%\usepackage[sc]{mathpazo}
%\usepackage{
%lmodern, 
%hfoldsty 
% charter
%}


%Language-related packages
\usepackage[english, portuguese,spanish]{babel} % defining languages 
\usepackage[babel]{csquotes} %English quotation marks
%\usepackage[babel, german=guillemets]{csquotes}

\usepackage[T1]{fontenc}
%\usepackage{lmodern}  % APA 7: reemplazado por Times (newtx)
% APA 7: tipografía Times New Roman 12 pt
\usepackage{newtxtext}
\usepackage{newtxmath}
\usepackage{eurosym}
% The € sign can now really be displayed with €
\DeclareUnicodeCharacter{20AC}{\euro}
\DeclareUnicodeCharacter{200B}{}
%The packages for the index
\usepackage{makeidx}
\makeindex

\usepackage{xcolor}
\usepackage{colortbl}

\usepackage{tikz}
\usetikzlibrary{calc,positioning}



%Various standard packages
\usepackage[
		pdftex,
		pdfpagelabels=true,
	 	citecolor=black,
       	filecolor=blue,
      	urlcolor=black,
       	bookmarks=true,
       	bookmarksopen=true,
       	bookmarksopenlevel=3,
       	plainpages=false,
       	pdfpagelabels=true,
       	pdfborder={0 0 0},
       	breaklinks=true,]{hyperref}
\hypersetup{colorlinks=true,
			breaklinks=true,
			urlcolor=black,
			linkcolor=black,
			menucolor=black,
			% Diese Angaben kommen in die PDF Eigenschaften
			pdftitle    = {Help-Seeking},
			pdfsubject  = {Dissertation},
			pdfauthor   = {Bastian Wimmer},
			pdfkeywords = {Keyword 1,Keyword 2, Keyword 3},
			pdfcreator  = {TeXStudio, MacTex \& Me},
			pdfproducer = {LaTeX with Hyperref}}
\usepackage[hypcap]{caption}



%%Line spacing
\usepackage{setspace}
%\setstretch{1,2381}
%\onehalfspacing
% ── FORMATO APA 7 ────────────────────────────────────────────
% Interlineado doble en todo el documento (texto, citas en bloque y referencias)
\doublespacing

% Sangría de primera línea de 0.5 pulgadas (1.27 cm) en todos los párrafos.
% El paquete indentfirst (cargado más abajo) la aplica también al primer
% párrafo que sigue a cada título, como exige APA.
\setlength{\parindent}{1.27cm}
% Espacio adicional entre párrafos para dar más aire a la lectura.
\setlength{\parskip}{6pt plus 2pt minus 1pt}

% Listas con el mismo interlineado doble del texto, sin espacio adicional
% entre elementos ni antes o después de la lista.
% (Se ajusta directamente porque enumitem es incompatible con paralist.)
\makeatletter
\def\@listi{\leftmargin\leftmargini
	\topsep\z@ \partopsep\z@ \parsep\z@ \itemsep\z@}
\let\@listI\@listi
\def\@listii{\leftmargin\leftmarginii \labelwidth\leftmarginii
	\advance\labelwidth-\labelsep
	\topsep\z@ \partopsep\z@ \parsep\z@ \itemsep\z@}
\makeatother

% Número de página en el extremo superior derecho, sin línea de encabezado,
% incluidas las páginas de inicio de capítulo (estilo plain).
\usepackage{fancyhdr}
\setlength{\headheight}{15pt}
\pagestyle{fancy}
\fancyhf{}
\fancyhead[R]{\thepage}
\renewcommand{\headrulewidth}{0pt}
\fancypagestyle{plain}{%
	\fancyhf{}%
	\fancyhead[R]{\thepage}%
	\renewcommand{\headrulewidth}{0pt}%
}

% Títulos de tablas y figuras según APA: el rótulo ("Tabla 1") en negrita,
% en su propia línea, y el título en cursiva debajo, alineados a la izquierda.
\captionsetup{labelfont=bf, textfont=it, labelsep=newline,
	justification=raggedright, singlelinecheck=false}

% ─────────────────────────────────────────────────────────────

% Packages for graphics, tables, etc.
\usepackage{verbatim} %for multiline comments
%\usepackage[nonumberlist, acronym, toc, section]{glossaries}
\usepackage{textcomp}
\usepackage{graphicx} % support the \includegraphics command and options
\usepackage{booktabs}
\usepackage{array} % for better arrays (eg matrices) in maths
\usepackage{subfig} % make it possible to include more than one captioned figure/table in a single float
\usepackage{multicol}
\usepackage{multirow}
\usepackage{ifthen}
\usepackage{float}
   %\usepackage{csquotes}
   
\usepackage{tabularx}
\usepackage{adjustbox} % Para controlar el ajuste de la tabla
\usepackage{ragged2e} % Para justificar texto en columnas p{}


\newcolumntype{C}[1]{>{\centering\arraybackslash}m{#1}} 
\newcolumntype{L}[1]{>{\raggedright\arraybackslash}m{#1}}
\newcolumntype{R}[1]{>{\raggedleft\arraybackslash}m{#1}}
\newcolumntype{Y}{>{\raggedright\arraybackslash}X} % Columna alineada izquierda
\usepackage{paralist}
\usepackage{siunitx}
\sisetup{
add-decimal-zero = false,
add-integer-zero = false,
}

%Command for the lines on the title page
\newcommand{\HRule}{\rule{\linewidth}{0.2mm}}
%Indent paragraph
%Sets the indentation depth of the first line after the heading
\usepackage{indentfirst}

% Improve horizontal lines in tables
\newcommand{\forloop}[5][1]{%
\setcounter{#2}{#3}%
\ifthenelse{#4}{#5\addtocounter{#2}{#1}%
\forloop[#1]{#2}{\value{#2}}{#4}{#5}}%
{}}
\newcounter{crcounter}
\newcommand{\compensaterule}[1]{%
\forloop{crcounter}{1}{\value{crcounter} < #1}%
{\vspace*{-\aboverulesep}\vspace*{-\belowrulesep}}}
\newcommand{\multirowbt}[3]{\multirow{#1}{#2}%
{\compensaterule{#1}#3}}



%Paquetes propios

\addto\captionsspanish{
	\renewcommand{\tablename}{Tabla}
}


\usepackage{titling}
\usepackage{changepage} % Cargar el paquete changepage
\usepackage{svg}


%Captions left
% (Desactivado: forzaba títulos en letra pequeña, contrario a APA 7; el formato
% de títulos de tablas y figuras se define arriba con \captionsetup.)
%\usepackage[font=footnotesize,labelfont=bf,singlelinecheck=false,format=plain,,justification=justified,indention=0cm]{caption}


%Colores para las tablas
% ── grayscale palette ────────────────────────────────────────
\definecolor{bubInc}{gray}{0.88}   % light  — Include
\definecolor{bubUnc}{gray}{0.58}   % mid    — Uncertain
\definecolor{bubExc}{gray}{0.20}   % dark   — Exclude
\definecolor{gridcol}{gray}{0.80}
\definecolor{textcol}{gray}{0.12}
\definecolor{subcol}{gray}{0.48}



%Citation in APA style
\usepackage[style=apa, backend=biber]{biblatex}  % Load biblatex with APA style
\addbibresource{./_bibliografia.bib}  % Ensure this matches your .bib file name
% APA 7: lista de referencias con sangría francesa (la aplica biblatex-apa) e
% interlineado doble, sin espacio adicional entre entradas.
\setlength{\bibitemsep}{0pt}

